% $Header: /home/vedranm/bitbucket/beamer/solutions/generic-talks/generic-ornate-15min-45min.en.tex,v 90e850259b8b 2007/01/28 20:48:30 tantau $

\documentclass[utf8x,t,aspectratio=169,xcolor={dvipsnames}]{beamer}

\mode<presentation>
{
%  \usetheme[height=7mm]{Rochester}
  \usetheme{Rochester}
  \useinnertheme{rounded}
%  \useoutertheme{infolines}
  % or ...

  \setbeamercovered{transparent}
}
\usefonttheme[onlymath]{serif}

\usepackage[english]{babel}
\usepackage[utf8x]{inputenc}

\usepackage{times}
\usepackage[T1]{fontenc}
% Or whatever. Note that the encoding and the font should match. If T1
% does not look nice, try deleting the line with the fontenc.

%% to draw trees without jpg figures
%\usepackage{synttree}
% For code blocks
% for sane tabular handling
%\usepackage{array}
% For strikethrough (normalem keeps it as emph)
%\usepackage[normalem]{ulem}

\usepackage{mathrsfs}
\usepackage{bm}  % For bold math

%\usepackage{marvosym}  % For tick
\usepackage{wasysym}  % For tick
\usepackage{adjustbox}

\newcommand{\vitem}{\item \vspace{4pt}}
\newcommand{\nyil}{$\rightarrow$\ }
\newcommand{\nagytilde}{$\sim$}
\newtheorem{Examplee}{Example}[section]
\providecommand{\abs}[1]{\lvert#1\rvert}
\providecommand{\norm}[1]{\lVert#1\rVert}

\newtheorem{nix}{}[section]

\title{Why implementation matters}
\subtitle{Evaluation of an open-source constraint grammar parser}

\author{Dávid Márk Nemeskey\inst{1}, Francis M. Tyers\inst{2}, Mans Hulden\inst{3}}
\institute{
  \inst{1} Institute for Computer Science and Control, Hungarian Academy of Sciences \\
  \inst{2} HSL-fakultetet, UiT Norgga árktalaš universitehta \\
  \inst{3} Department of Linguistics, University of Colorado Boulder}

\date % (optional)
{25th August 2014}

% If you have a file called "university-logo-filename.xxx", where xxx
% is a graphic format that can be processed by latex or pdflatex,
% resp., then you can add a logo as follows:

% \pgfdeclareimage[height=0.5cm]{university-logo}{university-logo-filename}
% \logo{\pgfuseimage{university-logo}}

\pgfdeclareimage[height=1.5cm]{apertium-logo}{apertium.png}
\logo{\pgfuseimage{apertium-logo}}

%\AtBeginSection[]
%{
%   \begin{frame}
%       \frametitle{Outline}
%       \tableofcontents[currentsection]
%   \end{frame}
%}
%

% If you wish to uncover everything in a step-wise fashion, uncomment
% the following command: 
%\beamerdefaultoverlayspecification{<+->}

\setbeamertemplate{navigation symbols}{}%remove navigation symbols

\begin{document}

% TODO: do we need an outline?
\begin{frame}{}
  \titlepage
\end{frame}

\begin{frame}{Outline}
\begin{itemize}
    \vitem Introduction
    \vitem Practice
    \begin{itemize}
        \vitem Optimisation
        \vitem Evaluation
    \end{itemize}
    \vitem Theory
    \begin{itemize}
        \vitem Complexity analysis
        \vitem The "$\log{X}$ experiment"
    \end{itemize}
    \vitem Conclusions
\end{itemize}
\end{frame}

\section{Introduction}

\begin{frame}{Constraint Grammar}
Rule-based formalism for linguistic tasks:
\begin{itemize}
    \vitem morphological analysis
    \vitem clause boundary detection
    \vitem surface syntactic parsing
\end{itemize}

Advantages of CG:
\begin{itemize}
    \vitem very high accuracy on free text
    \vitem works for languages without annotated corpora
\end{itemize}
\end{frame}

\begin{frame}{Constraint Grammar History}
Implementations:
\begin{itemize}
  \vitem CG: Introduced by Karlsson (1995)
  \vitem CG-2: Tapanainen (1996)
  \vitem VISL CG: Bick (1996)
  \vitem VISL CG3: Didriksen (2011)
\end{itemize}
Finite state experiments:
\begin{itemize}
    \vitem Hulden (2011)
    \vitem Peltonen (2011)
    \vitem Yli-Jyrä (2011)
\end{itemize}

%\textbf{Other relevant literature:}
%\begin{itemize}
%  \item Samuelsson, C., Voutilainen, A. (1997) ``Comparing a linguistic and a stochastic tagger''. \emph{35th ACL}. 246--253
%  \item ...
%\end{itemize}

\end{frame}

% What is CG?
% Apertium & VISL-CG; CG vs n-grame? 
\begin{frame}{Apertium}
Rule-based machine translation platform

\begin{itemize}
  \vitem Over 40 released language pairs
  \vitem Free/open-source: GPL licence
\end{itemize}
Many language pairs include constraint grammars
\begin{itemize}
  \vitem All pairs with constraint grammars use VISL CG-3
  \vitem Accuracy is higher than statistical ($n$-gram) models
  \vitem Main drawback: running time is much worse
\end{itemize} 

\end{frame}

\begin{frame}{Aims}
We set out with the following aims:

\smallskip

\begin{tabular}{ l  l }
  Measure actual speed of CG-based morphological disambiguation & \textbf{\Square} \\
  Improve on the performance of state-of-the-art open source finite-state CG & \textbf{\Square} \\
  Try to outperform VISL-CG & \textbf{\Square} \\
\end{tabular}

\bigskip

We have chose \emph{fomacg} (Hulden, 2011) as the base for our parser.
\begin{itemize}
    \vitem Open-source CG $\rightarrow$ FST compiler 
    \vitem Documented well in the paper
    \vitem All experiments used \emph{foma}; due to Mans Hulden
\end{itemize}
\end{frame}

\section{Methodology}

\begin{frame}{Methodology}
Development
\begin{itemize}
    \vitem Naïve parser implementation for baseline
    \vitem Optimisations
    \begin{itemize}
        \vitem parser
        \vitem compiler
    \end{itemize}
    \vitem Benchmarked on a small Hungarian CG grammar
\end{itemize}

Testing
\begin{itemize}
    \vitem The performance has been measured against VISL-CG.
    \vitem Three CGs: Hungarian, Breton, Finnish
\end{itemize}

\end{frame}

\begin{frame}{Development and test sets}
\begin{table}
  \centering
  \begin{tabular}{ | l | r | r | r | }
  \hline
  \textbf{Language} & \textbf{Rules} & \textbf{CG-2 rules} & \textbf{Corpus size} \\
  \hline
  Hungarian &   33 &   33 & 11 * 32 \\
  Breton    &  251 &  226 & 1,161 \\
  Finnish   & 1207 & 1172 & 1,620 \\
  \hline
  \end{tabular}
\end{table}
\end{frame}

\section{fomacg}

\begin{frame}{Morphological disambiguation based on CG}
Ambiguous input sentence (\emph{fomacg} format):

\uncover<2->{
  \begin{itemize}
    \uncover<2->{\vitem {\color<2>{red}cohorts}}
    \uncover<3>{\vitem {\color<3>{red}readings}}
  \end{itemize}
}

\bigskip

\texttt{\$0\$ "\textbf{All}" \#BOC\# | \#0\# "all" <det> | \#EOC\#}\\
\texttt{\$0\$ "\textbf{is}" \#BOC\# | \#0\# "be" <vbser> <pres> <s3p> | \#EOC\#}\\
\only<1>{\texttt{\$0\$ "\textbf{well}" \#BOC\# | \#0\# "well" <n> <sg> <nom> | \#0\# "well" <adj> <sg> <nom> | \#EOC\#}\\}
\only<2>{\texttt{\color{red}\$0\$ "\textbf{well}" \#BOC\# | \#0\# "well" <n> <sg> <nom> | \#0\# "well" <adj> <sg> <nom> | \#EOC\#}\\}
\only<3>{\texttt{\$0\$ "\textbf{well}" \#BOC\# | {\color{red}\#0\# "well" <n> <sg> <nom>} | {\color{red}\#0\# "well" <adj> <sg> <nom>} | \#EOC\#}\\}
\texttt{\$0\$ "\textbf{.}" \#BOC\# | \#0\# "." <punct> | \#EOC\#}

\end{frame}

\begin{frame}{Morphological disambiguation based on CG}
Rules (constraints):
\begin{itemize}
    \vitem select or filter readings
    \vitem grouped into priority sections
    \vitem two ways to look at it:
    \begin{itemize}
        \vitem constraint solving (optimality theory)
        \only<1>{\vitem iterative rule application}
        \only<2>{\vitem \textbf{iterative rule application}}
    \end{itemize}
    \vitem \emph{fomacg} converts these rules to FST
\end{itemize}

\bigskip

\texttt{REMOVE (N SG) IF (-1C VBSER)}
\end{frame}

\begin{frame}{Morphological disambiguation based on CG}
Rule application:

\bigskip
\bigskip
\bigskip

\texttt{\$0\$ "\textbf{All}" \#BOC\# | \#0\# "all" <det> | \#EOC\#}\\
\texttt{\$0\$ "\textbf{is}" \#BOC\# | \#0\# "be" <vbser> <pres> <s3p> | \#EOC\#}\\
\only<1>{\texttt{\$0\$ "\textbf{well}" \#BOC\# | \#0\# "well" <n> <sg> <nom> | \#0\# "well" <adj> <sg> <nom> | \#EOC\#}\\}
\only<2>{\texttt{\$0\$ "\textbf{well}" \#BOC\# | \textbf{\color{red}\#X\#} "well" <n> <sg> <nom> | \#0\# "well" <adj> <sg> <nom> | \#EOC\#}\\}
\texttt{\$0\$ "\textbf{.}" \#BOC\# | \#0\# "." <punct> | \#EOC\#}

\bigskip
\bigskip
\bigskip

\texttt{FST\{ REMOVE (N SG) IF (-1C VBSER) \}}
\end{frame}

\section{Generic optimisations}

\begin{frame}[fragile]
  \only<1>{\frametitle{CG parser algorithm}}
  \only<2-3>{\frametitle{Naïve baseline}}
\begin{tabular}{p{22em} p{8em}}
% TODO algorithm
\begin{semiverbatim}
for (sentence in text) \{
  fomacg\_str = convert\_in(sentence);
FIRST\_RULE:
  for (section in cg)
    for (rule in section)
      if rule.test\_and\_apply(fomacg\_str)
        goto FIRST\_RULE;
    print convert\_out(fomacg\_str);
\}
\end{semiverbatim}
\only<2-3>{
\begin{itemize}
    \vitem rules are FSTs
    \vitem convert functions are FSTs
\end{itemize}
}
&
\only<3>{
\smash{\raisebox{-\height}{\begin{tabular}{| l | r | r |}
  \hline
  \textbf{Method} & \textbf{Time} & \textbf{$\Delta$} \\
  \hline
  VISL-CG & 0.284s & -- \\
  \hline
  Naïve & 6.4s & 22.5x \\
  \hline
  \end{tabular}}}
}
\end{tabular}
\end{frame}

\begin{frame}{Discarded readings}
\begin{tabular}{p{22em} p{8em}}
The original rule FSTs keep discarded readings, only prepend them with \texttt{\#X\#}.

\bigskip

If we delete them:
\begin{itemize}
    \vitem after each rule application, the sentence gets shorter
    \vitem subsequent rules can be applied faster
\end{itemize}
&
\only<2>{
\smash{\raisebox{-\height}{\begin{tabular}{| l | r | r |}
  \hline
  \textbf{Method} & \textbf{Time} & \textbf{$\Delta$\%} \\
  \hline
  VISL-CG & 0.284s & -- \\
  \hline
  Naïve & 6.4s & -- \\
  DeleteX & 5.9s & 8\% \\
  \hline
  \end{tabular}}}
}
\end{tabular}

\end{frame}

\begin{frame}{Rule testing}
\begin{tabular}{p{22em} p{8em}}
Further speed-ups could be achieved by separating rule testing and application:
\begin{itemize}
    \vitem an FSA tests if the rule is applicable
    \vitem the rule FST is only applied if it is
\end{itemize}

\smallskip

Automata are faster than tranducers:
\begin{itemize}
    \vitem no output is generated
    \vitem can be determinised and minimised
\end{itemize}
\only<2>{After implementing it...}
&
\only<2>{
\smash{\raisebox{-\height}{\begin{tabular}{| l | r | r |}
  \hline
  \textbf{Method} & \textbf{Time} & \textbf{$\Delta$\%} \\
  \hline
  VISL-CG & 0.284s & -- \\
  \hline
  Naïve & 6.4s & -- \\
  DeleteX & 5.9s & 8\% \\
  Testing & 10s & -- \\
  \hline
  \end{tabular}}}
}
\end{tabular}
\end{frame}

\begin{frame}{Foma is too generic}

The \texttt{foma} library supports the following features:
\begin{itemize}
  \vitem Conversion of the text to symbols (single- and multi-character)
  \vitem Regular transitions and flag diacritics
  \vitem Three types of search in the transition matrix (linear, binary and indexed)
  \vitem Deterministic and non-deterministic operation
  \vitem Iterators (multiple invocations iterate the non-deterministic outputs)
  \vitem Transducer output
\end{itemize}
%When we run a FSA, all these features are taken into account.
\end{frame}

\begin{frame}{Cut down on features}

\setbeamercovered{transparent}
What we actually need...
\begin{itemize}
  \vitem Conversion of the text to symbols (single- and multi-character)
  \vitem Regular transitions \uncover<1>{and flag diacritics}
  \vitem \uncover<1-2>{Three types of} search in the transition matrix (\uncover<1-2>{linear, }binary\uncover<1-2>{ and indexed})
  \vitem Deterministic \uncover<1-3>{and non-deterministic }operation
  \vitem \uncover<1-4>{Iterators (multiple invocations iterate the non-deterministic outputs)}
  \vitem \uncover<1-5>{Transducer output}
\end{itemize}

\only<7>{Implementing our own functions may improve performance.}
\end{frame}

\begin{frame}{Performance of our custom functions}
\begin{tabular}{p{22em} p{8em}}
Deterministic FSA
\begin{itemize}
    \vitem rule testers
\end{itemize}
Deterministic FST
\begin{itemize}
    \vitem format conversion
\end{itemize}
Non-deterministic functional FST
\begin{itemize}
    \vitem rule transducers
\end{itemize}
&
\only<2>{
\smash{\raisebox{-\height}{\begin{tabular}{| l | r | r |}
  \hline
  \textbf{Method} & \textbf{Time} & \textbf{$\Delta$\%} \\
  \hline
  VISL-CG & 0.284s & -- \\
  \hline
  Naïve & 6.4s & -- \\
  DeleteX & 5.9s & 8\% \\
  Testing & 10s & -- \\
  Custom FSA & 1.45s & 75\% \\
  Custom DFST & 1.275s & 12\% \\
  Custom FFST & 1.05s & 18\% \\
  \hline
  \end{tabular}}}
}
\end{tabular}

\end{frame}

\begin{frame}{Learnings}
Genericity is useful:
\begin{itemize}
    \vitem it allows a library to be used in many situations  % TODO Ugh
\end{itemize}

When performance is important, relying on generic libraries
\begin{itemize}
    \vitem leads to suboptimal performance
    \vitem produces counterintuitive results
\end{itemize}

When doing research, it is important to use software the behaviour of which
reflects theory.

\end{frame}

\section{CG-specific optimisations}

\begin{frame}{CG-specific optimisations}
The changes thus were agnostic to the structure of the CG parsing task:
\begin{itemize}
    \vitem decreasing input length
    \vitem using automata for rule testing
    \vitem using custom functions
\end{itemize}

\only<2>{
The rule FSAs \& FSTs have special properties:
\begin{itemize}
    \vitem multicharacter symbols, separated by spaces 
    \vitem common symbols:
    \begin{itemize}
        \vitem boilerplate (\texttt{\$0\$}, \texttt{\#X\#}, etc.)
        \vitem POS, words (\texttt{N}, \texttt{the}, etc.)
    \end{itemize}
\end{itemize}
}
\end{frame}

\begin{frame}{String vector machines}
\begin{tabular}{p{22em} p{8em}}
Symbols are multicharacter, with spaces between them (e.g.
\texttt{\$0\$ "as" \#BOC\# | \#0\# "as" <PREP> | \#EOC\#})
\begin{itemize}
    \vitem Our FSTs are $char* \rightarrow char*$
    \vitem String to symbol vector transformation occurs on each FSA/FST application
\end{itemize}
Why not do it in a preprocessing step?
\begin{itemize}
    \vitem Do the transformation only once
    \vitem Rewrite the FSTs to $StringVect \rightarrow StringVect$
\end{itemize}
&
\only<2>{
\smash{\raisebox{-\height}{\begin{tabular}{| l | r | r |}
  \hline
  \textbf{Method} & \textbf{Time} & \textbf{$\Delta$\%} \\
  \hline
  VISL-CG & 0.284s & -- \\
  \hline
  Naïve & 6.4s & -- \\
  DeleteX & 5.9s & 8\% \\
  Testing & 10s & -- \\
  Custom FSA & 1.45s & 75\% \\
  Custom DFST & 1.275s & 12\% \\
  Custom FFST & 1.05s & 18\% \\
  StringVect & 1s & 4\% \\
  \hline
  \end{tabular}}}
}
\end{tabular}
\end{frame}

\begin{frame}{String to symbol conversion}
Internally foma uses integer symbol ids, so a $String \rightarrow int$ conversion
is still performed separately by each FSA/FSTs.
\begin{itemize}
    \vitem Usually taken for granted in the literature
    \vitem Significant contribution to execution time (60\%)
\end{itemize}
\only<2>{
foma uses per-FST tries to do the conversion
\begin{itemize}
    \vitem Doing a single conversion is fast, but it adds up
    \vitem Duplicated trie branches: boilerplate, CG tags
    \vitem High memory consumption
\end{itemize}
}
\end{frame}

\begin{frame}{Converter transducer}
\begin{tabular}{p{22em} p{8em}}
Do the conversion only once in a converter transducer $C$:
\begin{itemize}
    \vitem $\Sigma_c = \displaystyle \bigcup_{i=0}^N{\Sigma_{R_i}}$
    \vitem $C: String \rightarrow int$
\end{itemize}
Rest of the machines are simply $R_i: int \rightarrow int$.
\only<2>{
\begin{itemize}
    \vitem Running time improved by 40--70\%.
    \vitem Memory consumption reduced by 75--80\%.
\end{itemize}
}
&
\only<2>{
\smash{\raisebox{-\height}{\begin{tabular}{| l | r | r |}
  \hline
  \textbf{Method} & \textbf{Time} & \textbf{$\Delta$\%} \\
  \hline
  VISL-CG & 0.284s & -- \\
  \hline
  Naïve & 6.4s & -- \\
  DeleteX & 5.9s & 8\% \\
  Testing & 10s & -- \\
  Custom FSA & 1.45s & 75\% \\
  Custom DFST & 1.275s & 12\% \\
  Custom FFST & 1.05s & 18\% \\
  StringVect & 1s & 4\% \\
  One-time trie & 0.32s & 68\% \\
  \hline
  \end{tabular}}}
}
\end{tabular}
\end{frame}

\section{Complexity analysis}

\begin{frame}{Complexity analysis}
Tapanainen proved that worst-case time complexity for a CG parser is
$\mathcal{O}(n^3k^2G)$, where
\begin{itemize}
    \vitem $n$ is the length of the sentence;
    \vitem $k$ is the maximum number of readings per cohort;
    \vitem $G$ is the number of rules in the grammar.
\end{itemize}
In more detail,
\begin{itemize}
    \vitem testing a cohort with a rule is $\mathcal{O}(nk)$;
    \vitem the whole sentence, $\mathcal{O}(n^2k)$;
    \vitem for each rule: $\mathcal{O}(n^2kG)$;
    \vitem we need at most $n(k - 1)$ rounds.
\end{itemize}
\end{frame}

\begin{frame}{Finite-state complexity}
Hulden showed that with FST rules, complexity decreases to $O(n^{2}k^{2}G)$.
\begin{itemize}
    \vitem an FST applies the rule everywhere in the sentence in $\mathcal{O}(nk)$
    \vitem precisely, it takes $\mathcal{O}(nk \norm{FST})$,
           where $\norm{FST}$ can be large
    \vitem if the rules are factored to bimachines, this decreases $\mathcal{O}(2nk)$
    \vitem fomacg does not include bimachine factorisation
\end{itemize}
\only<2>{

\smallskip

Our work improved on these numbers by
\begin{itemize}
    \vitem introducing rule testing automata, which can be applied in $\mathcal{O}(nk)$
    \vitem implementing custom functions that can live up to the theory
\end{itemize}
}
\end{frame}

\begin{frame}{Bimachines vs rule condition FSAs}
On average, our work also outperforms the pure bimachine setup. Let
\begin{itemize}
    \vitem $G_a$: number of rules that can be applied to the sentence
    \vitem $G_u$: number of rules that cannot
    \vitem $G = G_a + G_u$
\end{itemize}
Then, $G_a$ rounds (each $\frac{G_u}{2}$ unsuccessful test and 1 application) take
\begin{itemize}
    \vitem only bimachines: $~2G_a(1 + \frac{G_u}{2})$
    \vitem condition FSAs: $~G_a(1 + 2 + \frac{G_u}{2})$
\end{itemize}
Since $G_a << G$ and $3 << G_u$, condition FSAs perform better.
\end{frame}

\begin{frame}{Rule testing vs application}
$G_a << G$, but by how much?

\bigskip

\begin{tabular}{| l | l | l | l |}
\hline
\textbf{Language} & \textbf{Testing} & \textbf{Application} & \textbf{Other} \\
\hline
Hungarian & 112,416 / 45.84\% & 2,560 / 7.20\% & 13,760 / 22.75\% \\
Breton & 612,006 / 62.40\% & 4,707 / 2.74\% & 11,921 / 4.41\% \\
Finnish & 6,290,609 / 87.17\% & $\sim$0\% & 81,017 / 3.57\% \\
\hline
\end{tabular}

\bigskip

Is there a way to decrease the number of rule checks?

\end{frame}

\begin{frame}{Hierarchical rule checking}
We use the fact that regular languages are closed under the union operator.
\begin{itemize}
    \vitem Let $\mathrm{FSA}_{G_a}$ and $\mathrm{FSA}_{G_b}$ be rule testing FSAs
           for rules $G_a$ and $G_b$
    \vitem Then, $\mathrm{FSA}_{G_{ab}} = \mathrm{FSA}_{G_a} \cup \mathrm{FSA}_{G_b}$
           is a rule testing FSA for $G_{ab}$ = $G_a \cup G_b$
    \vitem If $\mathrm{FSA}_{G_{ab}}$ is minimized, then it checks $G_{ab}$ in
           $\mathcal{O}(nk)$.
\end{itemize}

\smallskip

This allows for hierarchical rule checking:
\begin{itemize}
    \vitem First check if $\mathrm{FSA}_{G_{ab}}$ accepts a sentence
    \begin{itemize}
        \vitem If yes, check with $\mathrm{FSA}_{G_a}$
        \vitem If not, we saved a check!
    \end{itemize}
\end{itemize}
\end{frame}

\begin{frame}{The $\log G$ experiment}
For $G$ rules, we can build a tree with $\log G$ levels, reducing
\begin{itemize}
    \vitem the number of checks to $\log G$;
    \vitem the full running time to $\mathcal{O}(n^2k^2\log G)$.
\end{itemize}
% So what's the problem.
\only<2>{
Unfortunately, this does not work in practice.
\begin{itemize}
    \vitem the size of the unified automata grows exponentially
    \vitem this effects both memory consumption and binary size
    \vitem a single tree is infeasible
    \vitem compromise solution: a forest of 2-4 level trees.
\end{itemize}
}
\end{frame}

\begin{frame}{Comparison with VISL-CG}
% TODO center horizontally without pushing it down (table[h] and center don't work)
  \begin{tabular}{ | l | r | r | r | r | r | }
  \hline
  \textbf{Language} & \textbf{Algorithm} & \textbf{Initialisation} &
  \textbf{Disambiguation} & \textbf{Memory \%} & \textbf{File size} \\
  \hline
  Hungarian & VISL-CG & -- & 0.284s & -- & 9kB \\
  Hungarian & flat & 0.028s & 0.32s & 0.1\% & 60kB \\
  Hungarian & forest & 0.62s & \textbf{\color{OliveGreen}0.234s} & 1.5\% & 5.9MB \\
  \hline
  Breton    & VISL-CG & -- & \textbf{0.77s} & -- & 36kB \\
  Breton    & flat & 0.5s & 1.55s & 1.3\% & 1.5MB \\
  Breton    & forest & \textbf{\color{red}11.14s} & 1.05s & 4.4\% & 60MB \\
  \hline
  Finnish   & VISL-CG & -- & \textbf{1.78s} & -- & 184kB \\
  Finnish   & flat & 1.5s & 22.87s & 4.1\% & 7.2MB \\
  Finnish   & forest & \textbf{\color{red}20.75s} & 9.95s & 8.9\% & 198MB \\
  \hline
  \end{tabular}
\end{frame}

\section{Conclusions}

\begin{frame}{Conclusions}
We set out with the following aims:
\begin{tabular}{ l  l }
  Measure actual speed of CG-based morphological disambiguation & \textbf{\CheckedBox} \\
  Improve on the performance of state-of-the-art open source finite-state CG & \textbf{\CheckedBox} \\
  Try to outperform VISL-CG & \textbf{\Square} \\
\end{tabular}

\smallskip

\begin{itemize}
    \vitem Experimented with decreasing the time of applying a CG to
           $\mathcal{O}(n^2k^2\log G$; our method proved infeasible in practice
    \vitem Our main finding is that implementation matters
    \begin{itemize}
        \vitem genericity hurts performance
        \vitem it might prevent us from drawing valid conclusions from our
               experiments.
    \end{itemize}
\end{itemize}
\end{frame}

\begin{frame}{Acknowledgements}

\bigskip
\bigskip

This research was supported through an Apertium project in the 2013 Google Summer of Code.

\bigskip

\begin{table}
  \centering
  \begin{tabular}{ l  l }
    \adjustbox{valign=c}{\includegraphics[width=67.5pt]{apertium.png}}\hspace{2em} & \hspace{2em}\adjustbox{valign=c}{\includegraphics[width=67.5pt]{soc-logo-google-blue.jpg}} \\
  \end{tabular}
\end{table}
\end{frame}

\begin{frame}

\bigskip
\bigskip

\begin{center}
{\huge Thank you for your attention!}\\[\baselineskip]
{\huge Go raibh maith agaibh!}
\end{center}

\bigskip
\bigskip
\bigskip

emails: \\ {\tt nemeskey.david@gmail.com \\ francis.tyers@uit.no \\ mans.hulden@colorado.edu} \\

\end{frame}

\end{document}


